\clearpage
\VerdierPhysicalPageAnchor{0203}%
\begingroup\small\setlength{\parskip}{0pt}\setlength{\emergencystretch}{1.5em}
\setlength{\abovedisplayskip}{4pt plus 2pt minus 2pt}%
\setlength{\belowdisplayskip}{4pt plus 2pt minus 2pt}%

\VerdierSemanticAnchor{ch3d-2-4-1-cont}
\noindent\textit{是范畴等价。自然函子}
\[
K^{-}(M^{\wedge})/\operatorname{Ac}^{-}(M^{\wedge})
\longrightarrow D^{-}(\mathcal A)
\]
\noindent\textit{是充分忠实的。$D^{-}(\mathcal A)$ 中上同调对象均属于
$M^{\wedge}$ 的对象，同构于
$K^{-}(M^{\wedge})/\operatorname{Ac}^{-}(M^{\wedge})$ 的对象；这里
$\operatorname{Ac}^{-}(M)$ 与 $\operatorname{Ac}^{-}(M^{\wedge})$ 分别表示
分量属于 $M$ 与 $M^{\wedge}$ 的无环复形所成的范畴。}

\VerdierSemanticAnchor{ch3d-2-4-1-proof}
本命题立即由（2.3.1）、（2.3.2）及（第二章，2.3.5）推出。

\VerdierSemanticAnchor{ch3d-2-4-2}
\medskip\noindent\textbf{推论 2.4.2.}\quad
\textit{设 $U$ 是一个宇宙。假设 $\mathcal A$ 是一个 $U$-范畴，即对
$\mathcal A$ 的每一对对象 $Y$ 与 $Z$，
$\operatorname{Hom}_{\mathcal A}(Y,Z)$ 都是 $U$ 的元素，并且由 $U$ 的元素
作指标的直和在 $\mathcal A$ 中均可表示。设 $(X_i)_{i\in I\in U}$ 是
$\mathcal A$ 的一个生成元族；令 $M$ 为 $\mathcal A$ 中所有形如
$\coprod_{j\in J\in U}X_j$ 的对象所成的集合，其中 $X_j$ 属于集合
$\bigcup_{i\in I}\{X_i\}$。则自然函子}
\[
K^{-}(M)/\operatorname{Ac}^{-}(M)\longrightarrow D^{-}(\mathcal A)
\]
\noindent\textit{是范畴等价。}

\VerdierSemanticAnchor{ch3d-2-4-2-proof}
这立即由（2.4.1）推出。

\VerdierSemanticAnchor{ch3d-2-4-3}
\medskip\noindent\textbf{命题 2.4.3.}\quad
\textit{设 $\mathcal A$ 是阿贝尔范畴，$M$ 是 $\mathcal A$ 的对象集合，
具有性质 $(E_4)$ 与 $(E_5)$；假设 $\mathcal A$ 的每个对象都是 $M$ 中某个
对象的商对象，并且存在整数 $n\in\mathbb N$，使它成为 $\mathcal A$ 中每个
对象 $X$ 的 $M$-维数的上界。}

\noindent\textit{a) 范畴 $K^{*}(M)$ 关于
$\operatorname{Ac}^{*}(\mathcal A)$ 是左局部化的
（$*=-,+,b,{}$ ``空''）。}

\noindent\textit{b) 自然函子}
\[
K^{*}(M)/\operatorname{Ac}^{*}(M)\longrightarrow D^{*}(\mathcal A)
\]
\noindent\textit{是范畴等价。}

\VerdierSemanticAnchor{ch3d-2-4-3-proof}
这是命题（2.3.3）（$*=+,b,{}$ ``空''）或命题（2.3.1）（$*=-$）的改述。

\VerdierSemanticAnchor{ch3d-2-4-4}
\medskip\noindent\textbf{命题 2.4.4.}\quad
\textit{以
$Q^{*}\colon K^{*}(\mathcal A)\longrightarrow D^{*}(\mathcal A)$ 表示典范函子。}

\noindent\textit{a) 设 $X$ 是 $\mathcal A$ 的对象。对象 $X$ 是
$Q$-左自由的（第二章，2.3.2），当且仅当 $X$ 是范畴 $\mathcal A$ 的投射对象。}
\endgroup

